%%%%%%%%%%%%%%%%%%%%%%%%%%%%%%%%%%%%%%%%%%%%%%%%%%%%%%%%%%%%%%%%%%%%%%
% Lớp 11 | Dãy số. Cấp số cộng. Cấp số nhân | Cấp số cộng
% Dạng 1D2?2-1: Câu hỏi lý thuyết
% Số câu: 6  (sinh tự động bởi tach_cau.py; từ nay thầy sửa trực tiếp file này)
%%%%%%%%%%%%%%%%%%%%%%%%%%%%%%%%%%%%%%%%%%%%%%%%%%%%%%%%%%%%%%%%%%%%%%

%%%%%%%%%%%%%%%% TRẮC NGHIỆM 4 PHƯƠNG ÁN %%%%%%%%%%%%%%%%

% nguồn: Ok_TracNghiemLop11_Lan1
\begin{ex}%[1D2N2-1]%[ma:f4b0470a]
	Dãy số nào sau đây  là cấp số cộng?
	 \choice
		{\True $3$; $5$;  $7$; $9$; $11$}
		{$3$; $5$; $7$; $9$; $12$} 
		 {$3$; $5$; $7$; $9$; $10$}
		{$3$; $5$;  $6$; $7$; $9$}
	 \loigiai{
		 Xét dãy số $3$; $5$; $7$; $9$; $11$ có
		\begin{itemize}
			\item $5=3  + 2$.
			\item $7=5 +  2$.
			\item $9=7 + 2$. 
			\item $11=9 +  2$.
		\end{itemize}
		Đây là  cấp số cộng với $u_1=3$ công sai $d=2$.
	 }
\end{ex}

% nguồn: Ok_TracNghiemLop11_Lan1
\begin{ex}%[1D2N2-1]%[ma:bbabff3c]
	Cho  một cấp số cộng có $u_1=-3;u_2=3$. Tìm $d$.
	 \choice
		 {$d=5$}
		{$d=7$} 
		{\True $d=6$} 
		{$d=0$} 
	 \loigiai{
		Ta có $d=u_2-u_1=3+3=6$. 
	} 
\end{ex}

% nguồn: Ok_TracNghiemLop11_Lan1
\begin{ex}%[1D2N2-1]%[ma:62611b2f]
	Cấp số cộng với số hạng đầu là $u_1$, công sai $d$ có  số hạng tổng quát là
	 \choice
		{$u_n  = u_1 - (n-1)d, (n\geq 2)$}
		{\True $u_n =  u_1 + (n-1)d, (n\geq 2)$}
		{$u_n = u_1 + d, (n\geq 2)$} 
		{$u_n = u_1 + (n+1)d, (n\geq  2)$}
	 \loigiai{
		Số hạng tổng quát là $u_n =  u_1 + (n-1)d, (n\geq 2)$.
	} 
\end{ex}

% nguồn: Ok_TracNghiemLop11_Lan1
\begin{ex}%[1D2H2-1]%[ma:e728a8cb]
	Cho cấp số cộng có số hạng đầu tiên là $u_1 = 17$,  $u_5 = 37$. Công sai của cấp số cộng là
	 \choice
		{$10$} 
		 {$2$}
		{$8$} 
		{\True  $5$}
	 \loigiai{
		Ta  có $$u_5 = u_1 + (5-1)d \Rightarrow d = \dfrac{u_5 - u_1}{5 - 1} = 5.$$
		Vậy công sai của cấp số cộng này là $d = 5$. 
	 }
\end{ex}


%%%%%%%%%%%%%%%% ĐÚNG/SAI %%%%%%%%%%%%%%%%

% nguồn: Ok_dungsai_lop11_Lan1
\begin{ex}%[1D2H2-1]%[ma:af625934]
 Cho cấp số cộng $(u_n)$  có số hạng đầu $u_1$ và công sai $d$, xét tính đúng sai các mệnh đề sau
  \choiceTF
  {\True Số hạng tổng quát $u_n=u_1+(n-1)d$, với  $n\ge 1$}
  {Số  hạng $u_{n+1}=u_n+(n-1)d$, với $n\ge 1$}
  {\True Tổng  của $n$ số hạng đầu tiên của cấp số cộng là $S_n=\dfrac{n(u_1+u_n)}{2}$}
   {\True Số hạng tổng quát $u_{n}=\dfrac{u_{n-1}+u_{n+1}}{2}$ với $n\ge 2$}
  \loigiai{
  \begin{itemchoice}
   \itemch Số hạng tổng quát  của cấp số cộng $u_n=u_1+(n-1)d$, với $n\ge 1$.\\
   Vậy mệnh đề ``Số hạng tổng quát $u_n=u_1+(n-1)d$, với $n\ge 1$''  là mệnh đề đúng.
   \itemch Theo  định nghĩa cấp số cộng thì $u_{n+1}=u_n+d$ với $n\ge 1$.\\
   Vậy mệnh đề ``Số hạng $u_{n+1}=u_n+(n-1)d$, với $n\ge 1$'' là  mệnh đề sai.
   \itemch Tổng của $n$ số hạng đầu tiên của cấp  số cộng là $S_n=\dfrac{n(u_1+u_n)}{2}$.\\
   Vậy mệnh đề ``Tổng của $n$ số hạng đầu  tiên của cấp số cộng là $S_n=\dfrac{n(u_1+u_n)}{2}$'' là mệnh đề đúng.
   \itemch Ta có mỗi số hạng (trừ số hạng cuối  đối với cố số cộng hữu hạn) đều là trung bình cộng của hai số hạng đứng kề nó trong dãy.\\
    Vậy mệnh đề ``Số hạng tổng quát $u_n=\dfrac{u_{n-1}+u_{n+1}}{2}$ với $n\ge 2$'' là mệnh đề đúng.
  \end{itemchoice}
  }
\end{ex}


%%%%%%%%%%%%%%%% TRẢ LỜI NGẮN %%%%%%%%%%%%%%%%

% nguồn: Ok_TraLoiNgan_Lop11_Lan1
\begin{ex}%[1D2T2-1]%[ma:a64a0859]
	Người ta trồng $465$ cây trong một khu vườn hình tam giác theo cách  sau: Hàng thứ nhất có $1$ cây, hàng thứ hai có $2$ cây, và cứ như thế mỗi hàng sau sẽ có nhiều hơn hàng ngay trước đó $1$ cây. Hỏi trong khu vườn có bao nhiêu hàng cây?
	\shortans{30} 
	 \loigiai{
		Số cây trên  mỗi hàng lập thành một cấp số cộng $\left(u_n\right)$ với số hạng đầu $u_1=1$ và công sai $d=1$.\\
		Gọi $n$ là số hàng cây trồng trong vườn, do đó $465$  là tổng số cây của $n$ hàng cây.\\
		 Ta có
		\allowdisplaybreaks 
		\begin{eqnarray*}
			S_n=465 &\Leftrightarrow& \dfrac{n(n+1)}{2}=465 \Leftrightarrow n^2+n-930=0\\ 
			&\Leftrightarrow&  \hoac{& n=30 &\quad& \text{(nhận)} \\ & n =-31. &\quad& \text{(loại)}}
		\end{eqnarray*}
		Vậy số hàng cây trong khu  vườn là $30$ hàng.
	 }
\end{ex}
